
\documentclass{report}

\input{preamble}
\input{macros}
\input{letterfonts}

\title{\Huge{Diritto}\\Sistemi Informativi Aziendali}
\author{\huge{Caberlotto Giovanni}}
\date{}

\begin{document}

\maketitle
\newpage% or \cleardoublepage
% \pdfbookmark[<level>]{<title>}{<dest>}
\pdfbookmark[section]{\contentsname}{toc}
\tableofcontents
\pagebreak

\chapter{Lo Stato e la Costituzione}
\section{La Costituzione Italiana: nascità, caratteri e struttura}

\subsection{Definizione di costituzione}

\dfn{Costituzione}{Una costituzione è la legge fondamentale di uno stato. Essa contiene dei principi fondamentali che sono alla base dell'ordinamento giuridico di uno stato. Negli stati contemporanei si presenta sotto forma di documento scritto in quanto costituisce l'ordinamento giuridico, il cui scopo è quello di proclamare i diritti inviolabili dei cittadini e di porre un limite ai poteri dello Stato.}
\nt{Per quanto riguarda gli stati moderni fa eccezzione la Gran Bretagna che non ha una costituzione scritta, ma le norme costituzionali rappresentano norme consuetudinarie, in quanto emanate in periodi storici molto differenti}

\subsection{Nascità della Costituzione italiana}
Per comprendere al meglio come è nata la Costituzione italiana occorre conoscere i fatti storici dalla concessione dello Statuo Albertino fino all'entrata in vigore della costituzione.
La prima data che incontriamo è il 1848 qunado Re Carlo Alberto di Savoia concesse al popolo del regno d'italia una carta costituzionale chiamata appunto Statuto Albertino.
Lo Statuto Albertino era una carta costituzionale che presentava caratteristiche diverse rispetto all'attuale carta costituzionale italiana, innanzitutto era flessibile il ciò comportava che fosse facilmente modificabile.
Non fu votato ma fu concesso, non fù il popolo a votare per l'entrata in vigore di esso e non fù il popolo a deciderne il contenuto. 
In fine era breve perchè si limitava a descrivere 84 articoli che regolavano il funzionamento dello stato e le principali libertà individuali.
Seguendo l'ordine cronologico si passa all'unità d'italia avvenuta nel 1861 con Re Vittorio Emanuele II in cui l'italia fù uno stato liberale monarchico con una carta costituzionale rappresentata dallo Statuto Albertino.
Da li a poco tutta l'Europa, Italia compresa, passò un periodo difficile causato dalla prima guerra mondiale, tuttavia il periodo che interessa a noi per comprendere al meglio come è nata la costituzione è l'italia del 1918, quando l'italia fu in preda a una crisi profonda che portò a grandi tensioni sociali, la borghesia industriale e terriera chiese allo stato di attuare una politica autoritaria che ponesse fine alle lotte popolari.
A interessarsi della questione fù il PNF (Partito nazionale fascista), guidato da Benito Mussolini, al quale nel 1922 re Vittorio Emanuele Affidò l'incarico di capo del governo. 
In pochi anni il PNF impose delle riforme che portarono l'italia ad essere uno stato dittatoriale, questo fu dovuto anche dalla flessibilità dello Statuto Albertino, il quale restò in vigore durante tutto il periodo fascista.
Il PNF trasformò in maniera radicale anche le libertà e i diritti dei cittadini:
\begin{itemize}
  \item Furono abolite le libertà civili dei cittadini e i partiti d'opposizione;
  \item Venne anche attuata una politica reppressiva verso le opposizioni al regime, minoranze linguistiche e religiose che portarono alla emanazione delle leggi razziali del 1938.
\end{itemize}

Con la fine della seconda guerra mondiale in particolare nel 1943 il fascimo in Italia cessò di esistere, difatti Re Vittorio Emanuele III con i poteri a lui attribuiti dall'ancora in vigore Statuto Albertino revoco Benito Mussolini leader del PNF dal titolo di capo di governo.
Tuttavia l'italia del post fascimo pur avendo riconquistato le libertà violate in epoca fascista ebbe due problematiche da affrontare: 
\begin{itemize}
  \item La questione istituzionale;
  \item L'elaborazione di un nuovo testo costituzionale.
\end{itemize}
Si arriva così al 2 giugno 1946 quando si svolsero le prime votazioni a suffragio effettivamente universale della storia italiana, lo scopo di questa votazione fu la risoluzione delle due problematiche citate precedentemente. L'Italia passo ufficialmente da monarchia a repubblica e inoltre fù eletta l'assemblea costituente la quale darà vita alla nuova carta costituzionale.
I partiti padri della costituzione furono la democrazia cristiana il partito comunista e il partito socialista, il successo della costituzione fu dovuto dal fatto che la costituzione nacque da un compromesso tra gli orientamenti politici nell'assemblea costituente senza che essi si presero vantaggi momentanei dall'emanazione di essa, la costituzione difatti fu emanata con uno sguardo verso il futuro.

\subsection{Caratteri e struttura della costituzione}

La Costituzione è la legge fondamentale da cui discendono e a cui si ispirano tutte le norme ordinarie. Nessuna legge difatti può entrare in contrasto con quanto dice la costituzione e, nel caso in cui ciò avvenga è sempre la Costituzione a prevalere.
Si parla quindi di Costituzione rigida, ciò significa che le norme costituzionali non possono essere cambiate da leggi ordinarie. Per modificare la Costituzione sono necessarie leggi particolari, le leggi di revisione costituzionale, per l'approvazione delle quali è prevista una procedura complessa che mira ad allontanare il rischio di facili cambiamenti (art. 138 della Costituzione),
\nt{Le leggi di revisione costituzionale per essere approvate prevedono una doppia approvazione delle camere all'interno del parlamento}
rigida quindi, non significa immodificabile; la procedura di revisione prevede però regole più complesse rispetto a quelle necessarie per approvare le leggi ordinarie, proprio perché, i cambiamenti da apportare al testo costituzionale devono essere frutto di un accordo che coinvolga tutte le forze politiche, non un'imposizione della volontà della maggioranza. \nt{Inoltre l'articolo 139 della Costituzione, specifica che la forma reppubblicana non può essere oggetto di revisione costituzionale.}
La nostra Costituzione è lunga perché non si limita a sanciri i principi fondamentali, ma riconosce una pluralità di diritti anche riguardo ai rapporti sociali, etici ed economici. é infatti formata da 139 articoli, essa si apre con una parte introduttiva "I principi fondamentali" (art. 1-12 Costituzione), poi seguono la prima parte (art. 13-54 Costituzione) intitolata diritti e doveri dei cittadini, la quale tratta del rapporto tra i cittadini, segue poi la seconda parte, (art 55-139 Costituzione), intitolata ordinamento della repubblica, che tratta ell’organizzazione dei pubblici
poteri;seguono le disposizioni transitorie e finali (18 articoli) contenenti in prevalenza norme transitorie per il
passaggio al nuovo ordinamento.
Riassumendo la costituzione italiana è: 
\begin{itemize}
  \item Rigida;
  \item Lunga;
  \item Votata;
  \item Nata da un compromesso (contratto costituzionale).
\end{itemize}

è divisa in tre parti:

\begin{itemize}
  \item I principi fondamentali (artt. 1-12);
  \item Prima parte "Diritti e doveri dei cittadini"(artt. 13-54) ;
  \item Seconda parte "Ordinamento della repubblica" (artt. 55-139);
  \item Infine vi sono le disposizioni transitive e finali 18 articoli transitori per il passaggio al nuovo ordinamento.
\end{itemize}

\section{Principi fondamenti della Costituzione italiana}

I principi fondamentali (artt.1-12) sono una specie di “preambolo”, nel quale sono delineate le linee portanti della Carta
costituzionale: il principio democratico, lavorista, solidarista, di uguaglianza, di libertà.
Questi principi rappresentano i presupposti della società e gli obiettivi verso i quali essa deve tendere.
In questi primi articoli della Carta costituzionale si ritrovano gli ideali della rivoluzione francese, sintetizzati
egregiamente nell’articolo 1 (la libertà), nell’articolo 2 (la fraternità, o la solidarietà sociale) e nell’articolo 3
(l’uguaglianza).
Più in particolare, in questa “premessa” si delineano:

\begin{itemize}
  \item La forma di Stato (art.1);
  \item I rapporti tra lo Stato, i cittadini e i diversi soggetti di diritto (artt. 2, 3 , 4, 6);
  \item I rapporti tra lo Stato e le autonomie locali (art. 5);
  \item I rapporti tra lo Stato, la chiesa cattolica e le altre confessioni religiose (artt. 7 e 8);
  \item La promozione della cultura e la tutela del paesaggio (art. 9);
  \item I rapporti tra o Stato e l’ordinamento internazionale (artt. 10 e 11);
  \item Il simbolo dello Stato stesso (art. 12).
\end{itemize}

\subsection{Art 1 - Il principio democratico e lavorista}
\textit{“L’Italia è una Repubblica democratica fondata sul lavoro.
La sovranità appartiene al popolo che la esercita nelle forme e nei limiti della costituzione."}
\\
Il comma 1 sancisce il concetto di "Repubblica democratica" (sono l’una il rafforzativo dell’altra) cioè la forma di
governo nella quale tutte le cariche pubbliche si riconducono direttamente o indirettamente al consenso del popolo.
Viene inoltre sancito il principio del lavoro.
Il comma 2 sancisce il concetto di "Sovranità", si intende il potere supremo di governo, che appartiene esclusivamente
al popolo nella sua globalità, ma esercitabile solamente nei modi e nelle forme previste dalla Costituzione (art. 48 e
seguenti).

E' importante comprendere il modo nel quale i cittadini possono esercitare il potere di sovranità, vi sono due tipi di democrazia col quale un cittadino può esercitare il potere di sovranità: 
\begin{itemize}
  \item Democrazia indiretta;
  \item Democrazia diretta;
\end{itemize}

\textbf{Democrazia indiretta:}
Il popolo elegge i rappresentanti votando.
\textbf{Democrazia diretta:}
Vi sono tre modi col quale si esercita la democrazia diretta:
\begin{itemize}
  \item Con un referendum, il quale può essere consultivo o approvativo, con il referendum abbrogatorio inoltre il popolo può decidere se abolire una determinata legge, fanno eccezzione le leggi tributarie e di bilancio come enunciato dall'articolo 75.
  \item Petizione con la quale i cittadini possono esporre al parlamento necessità richieste ecc. 
  \item Richieste di Iniziative legislative, con cui 50.000 cittadini possono esporre al parlamento un progetto di legge purché sia esaminato e in caso approvato dal parlamento
\end{itemize}

\subsection{Art 2 -  Diritti inviolabili dell'uomo}
\textit{"La Repubblica riconosce e garantisce i diritti inviolabili dell’uomo, sia come singolo sia nelle formazioni sociali
dove si svolge la sua personalità, e richiede l’adempimento dei doveri inderogabili di solidarietà politica, economica e
sociale."}

\\

\nt{Principio di personalità}
Quest’articolo enuncia che lo Stato italiano, comprese le regioni, province, comuni e tutti coloro che si trovano nel
territorio italiano, riconosce quei diritti inviolabili (elencati nella parte prima dei Diritti e Doveri dei cittadini), vale a
dire quei diritti riconosciuti a tutti che nessuna legge può infrangere, cioè sopprimere e che nessuna persona può violare;
lo Stato non solo deve rispettarli ma deve anche proteggerli dalle violazioni provenienti dai soggetti privati. Con la
parola “riconosce” si intende che i diritti inviolabili fanno parte del patrimonio d’ogni individuo, sia come singolo
(diritto al nome, alla libertà d’espressione, ecc), sia come membro d’organizzazioni sociali (famiglia, partiti politici,
ecc). 

\nt{Principio di solidarietà}
Inoltre è introdotto il principio di solidarietà politica, economica e sociale. I doveri di solidarietà politica si
riferiscono a situazioni in cui la persona è chiamata a partecipare alla vita della comunità di cui fa parte (diritto di voto).
Adempiere ai doveri di solidarietà economica significa agire non pensando solo al nostro tornaconto, ma considerare
anche gli altri (pagare le tasse). Adempiere ai doveri di solidarietà sociale significa mettersi a disposizione gratuita di
chi ha bisogno (volontariato).

\subsection{Art 3 - L'uguaglianza}
\textit{“Tutti i cittadini hanno pari dignità sociale e sono eguali davanti alla legge, senza distinzioni di sesso, di
razza, di lingua, di religione, di opinioni politiche, di condizioni personali e sociali.
E’ compito della Repubblica rimuovere gli ostacoli di ordine economico e sociale, che, limitando di fatto la libertà e
l’uguaglianza dei cittadini, impediscono il pieno sviluppo della persona umana e l’effettiva partecipazione di tutti i
lavoratori all’organizzazione politica, economica e sociale del Paese.”}

\\

\dfn{Uguaglianza Formale}{Ciascun cittadino ha pari dignità di fronte alla legge.}

Nel c. 1 viene sancito il principio di uguaglianza formale di tutti davanti alla legge, questa espressione ha due
significati:

\begin{itemize}
  \item la legge è uguale per tutti (sia governanti che governati);
  \item tutti hanno gli stessi diritti, quindi non ci possono essere delle leggi che discriminano il sesso, la religione, la razza,
la lingua, l’opinione politica, le condizioni personali e sociali. 
\end{itemize}


\dfn{Uguaglianza Sostanziale}{Rappresenta la realizzazione dell'uguaglianza formale fra i cittadini.}

Nel c. 2 viene sancito il principio di uguaglianza sostanziale di tutti davanti alla legge proprio perché si riconosce che
nella realtà sono presenti delle discriminazioni, quindi la Repubblica deve essere in grado di offrire pari opportunità a
tutti affinché essi abbiano gli stessi diritti, rimuovendo qualsiasi tipo di ostacolo che impedisca ai cittadini di partecipare
alla vita politica, economica e sociale del Paese (abbattere barriere architettoniche, aiuti economici per le famiglie
bisognose, pensione sociale per tutti gli inabili al lavoro).

\ex{Aiuti dallo stato per attuare il principio di uguaglianza possono essere ad esempio verso famiglie con isee basso}



\subsection{Art 4 - Il principio lavorista}

\textit{“La Repubblica riconosce a tutti i cittadini il diritto al lavoro e promuove le condizioni che rendono effettivo questo
diritto.
Ogni cittadino ha il dovere di svolgere, secondo le proprie possibilità e la propria scelta, un’attività o una funzione che
concorra al progresso materiale o spirituale della società.”}

\\

Nel comma 1 viene sancito il principio del diritto al lavoro ma con tale espressione non si intende che lo stato è
obbligato a trovare un lavoro a tutti. Si intende invece che lo stato deve favorire l’economia e l’ingresso nel mondo del
lavoro.

Nel comma 2 Il lavoro oltre che un diritto è anche un dovere morale. Si ribadisce che il lavoro deve essere onesto.
\nt{Il cittadino deve svolgere una attività lavorativa per uno sviluppo per lo sviluppo personale e della collettività}

\subsection{Art 5 - Unità nazionale e autonomie locali}

\textit{“La Repubblica, una e indivisibile, riconosce e promuove le autonomie locali; attua nei servizi che dipendono dallo
Stato il più ampio decentramento amministrativo; adegua i principi ed i metodi della sua legislazione alle esigenze
dell'autonomia e del decentramento.”}

Ciò significa che l’Italia non è una composizione di Stati indipendenti che manifestano la volontà di federarsi, il che
esclude già da qui la possibilità che il nostro Stato possa – tramite una normale legge del parlamento – diventare uno
Stato federale. E’ piuttosto uno Stato Unitario regionale che riconosce il principio del decentramento
amministrativo. Il decentramento amministrativo è il principio secondo cui lo Stato non agisce soltanto con
organi centrali, ma si articola in enti autonomi locali (come i Comuni, le Province, le Città metropolitane e, in
particolare le Regioni) ed esercita le sue funzioni amministrative attraverso organi e uffici periferici (per esempio, l’ex
provveditorato agli studi, con competenza provinciale, rispetto al Ministero della Pubblica istruzione). Nei maggiori
Comuni, il decentramento è stato attuato con l'istituzione dei consiglio circoscrizionali o di quartiere.

\subsection{Art 6 - La tutela delle minioranze linguistiche}

\textit{"La Repubblica tutela con apposite norme le minoranze linguistiche."}

L’articolo 6 è uno dei tanti casi in cui il testo costituzionale ribadisce l’impegno ad eliminare tutti gli ostacoli che
limitino l’uguaglianza de cittadini: attraverso questo articolo in particolare, la Costituzione prescrive l’obbligo di
tutelare le minoranze linguistiche.
I costituenti hanno così reagito alle discriminazioni che in passato, e soprattutto durante il regime fascista, furono
attuate contro coloro che parlavano una lingua diversa.
In Alto Adige e in Valle d’Aosta è stato riconosciuto il bilinguismo (tedesco – italiano, francese – italiano): qui il
cittadino ha il diritto di usare ufficialmente la propria lingua e la Pubblica amministrazione è tenuta a rispettare tale
diritto, dotando i propri uffici (pubblici, giustizia e scuola) di personale bilingue.
In seguito alla recente riforma del titolo V della Costituzione (L.cost.n.3/2001) il bilinguismo è entrato anche nella
Costituzione in quanto, proprio la Valle d’Aosta e il Trentino Alto Adige sono citati con il nome bilingue.
L’articolo 6 è stato attuato di recente riguardo ai gruppi minori come i greci, albanesi, sloveni, croati, friulani, franco
provenzali ecc. che vivono in piccole comunità in diverse zone d’Italia, così la cultura e la lingua di queste popolazioni
sono state finalmente tutelate.

\nt{Durante il fascismo venne imposta una politica repressiva verso le minoranze linguistiche, l'articolo 6 nasce dall'esigienza di evitare che ciò si ripeta}


\subsection{Art 7 - I rapporti tra lo stato e la chiesa cattolica}

\textit{“Lo Stato e la Chiesa cattolica sono, ciascuno nel proprio ordine, indipendenti e sovrani.
I loro rapporti sono regolati dai Patti Lateranensi. Le modificazioni dei Patti, accettate dalle due parti, non richiedono
procedimento di revisione costituzionale.”}

Ai tempi dello Statuto Albertino l’Italia era uno stato confessionale cioè aveva un’unica religione si Stato che era la
religione cattolica.
I Padri costituenti vollero invece che l’Italia fosse uno stato laico riconoscendo alla Chiesa il potere religioso.
Comunque il fatto che abbiano riservato alla Chiesa cattolica un articolo, ci fa capire che essa aveva un ruolo
di grande importanza perché la stragrande maggioranza degli italiani erano cattolici e perché Roma è sede del Papato.
Nel c.2 viene stabilito che i rapporti tra lo stato e la Chiesa sono regolati dai Patti Lateranensi. Questi furono stipulati
nel 1929 tra Mussolini e il Papa Pio .
Per superare le contraddizioni tra art.7 e patti Lateranensi, nel 1984 venne stipulato un nuovo Concordato che
prevedeva:

\begin{itemize}
  \item Insegnamento della religione cattolica in modo facoltativo;
  \item Matrimonio civile, religioso e concordatario.
\end{itemize}

\subsection{Art 8 - I rapporti tra lo stato e le altre confessioni religiose}

\textit{“Tutte le confessioni religiose sono ugualmente libere davanti alla legge.
Le confessioni religiose diverse dalla cattolica hanno diritto di organizzarsi secondo i propri statuti, in quanto non
contrastino con l’ordinamento giuridico italiano.
I loro rapporti con lo Stato sono regolati per legge sulla base di intese con le relative rappresentanze.”}

Nel c.1 viene sancito il principio del pluralismo religioso cioè il riconoscimento della libertà religiosa di
ciascuno.

Ma già nel c.2 viene sancito un primo limite all’esercizio delle religioni diverse dalla cattolica e cioè che non devono
andare contro le leggi dello Stato italiano (ad. Es. in Italia è proibita ogni forma di menomazione fisica per cui
non può essere praticata l’infibulazione).
Nel c.3 si stabilisce che se una religione diversa dalla cattolica vuole vedersi riconosciuti dei diritti, deve
stipulare degli accordi con lo Stato italiano e l’iniziativa degli stessi deve partire dai rappresentanti delle
religioni.

\subsection{Art 9 - Tutela della cultura della ricerca e del paesaggio}

\textit{“La Repubblica promuove lo sviluppo della cultura e la ricerca scientifica e tecnica.
Tutela il paesaggio e il patrimonio storico e artistico della Nazione.”}

La Costituzione garantisce la massima libertà nella formazione e diffusione della cultura e nello svolgimento
dell’attività di ricerca, in contrapposizione alla politica seguita dal fascismo che aveva imposto la cultura di
regime. Non dà però una definizione di cultura che va quindi rintracciata nell’art.33.
Il c.2 è di estrema importanza perché contiene un concetto, quello di paesaggio, che ha subito nel corso del tempo una
profonda evoluzione. In assemblea costituente con tale termine si indicava unicamente la conservazione delle
bellezze naturali secondo quanto stabilito da una legge del 1938.
Oggi invece si intende un concetto molto più ampio di tutela del paesaggio che oggi prende il nome di tutela
dell’ambiente. Secondo un principio dello sviluppo economico-sociale la rigenerazione delle risorse non deve
compromettere l’ambiente delle generazioni future; comunque la legislazione riguardante la tutela dell’ambiente
si è avuta solo in tempi recenti.
Nel 1966 è stata emanata la prima legge antismog che detta le norme dell’inquinamento e solo nel 1986 è stato istituito
il Ministero dell’ambiente; soprattutto a partire dal 2000 che sono stati recepiti nel nostro ordinamento
importanti provvedimenti comunitari per quanto riguarda l’inquinamento atmosferico.

\ex{Un esempio è rappresentato dagli incentivi per le visite ai musei o i finanziamente da parte dello stato per esempio verso la sanità}

\subsection{Art 10 - L'Italia e la comunità internazionale}

\textit{“L'ordinamento giuridico italiano si conforma alle norme del diritto internazionale generalmente riconosciute.
La condizione giuridica dello straniero è regolata dalla legge in conformità delle norme e dei trattati
internazionali.
Lo straniero, al quale sia impedito nel suo paese l'effettivo esercizio delle libertà democratiche garantite dalla
Costituzione italiana, ha diritto d'asilo nel territorio della Repubblica secondo le condizioni stabilite dalla
legge.
Non è ammessa l'estradizione dello straniero per reati politici.”}

\nt{L'Italia si considera parte di un ordinamento più vasto: l'ordinamento internazionale}

Nel primo comma con la dicitura “L'ordinamento giuridico italiano si conforma alle norme del diritto
internazionale generalmente riconosciute” i padri costituenti vollero intendere che le norme internazionali che hanno
una portata generale (norme consuetudinarie) valgono automaticamente all’interno dell’ordinamento
giuridico italiano. Con questa norma costituzionale, inoltre, lo stato italiano si impegna a non adottare leggi di contrasto
con le norme del diritto internazionale, a considerare incostituzionali quelle leggi che non rispettassero tale principio e a
ratificare i trattati internazionali.
Nel secondo comma viene determinata la condizione giuridica dello straniero che è regolata dalla legge in conformità
delle norme e dei trattati internazionali.
Nel nostro ordinamento esistono attualmente due categorie di stranieri: i cittadini dell'Unione europea che godono di
una tutela e di garanzie simili a quelle del cittadino italiano; i cittadini extracomunitari, non appartenenti
all'Unione europea, che possono essere soggetti a restrizioni per quanto riguarda l'ingresso e la permanenza nel nostro
paese.
Non esiste una legge specifica per la questione degli stranieri ma l’unica che stabilisce diritti e doveri dello straniero è
la Legge Bossi-Fini del 2002. Essa prevede che l'espulsione, emessa in via amministrativa dal Prefetto della Provincia
dove viene rintracciato lo straniero clandestino, sia immediatamente eseguita con l'accompagnamento alla
frontiera da parte della forza pubblica. Gli immigrati clandestini, privi di validi documenti di identità, vengono
portati in centri di permanenza temporanea, istituiti dalla Legge Turco-Napolitano, al fine di essere identificati.
La legge prevede il rilascio del permesso di soggiorno, della residenza e cittadinanza italiana alle persone che
dimostrino di avere un lavoro o un reddito sufficienti per il loro mantenimento economico. A questa regola generale si
aggiungono i permessi di soggiorno speciali e quelli in applicazione del diritto di asilo. La norma ammette i
respingimenti al Paese di origine in acque extraterritoriali, in base ad accordi bilaterali fra Italia e Paesi limitrofi,
che impegnano le polizie dei rispettivi Paesi a cooperare per la prevenzione dell'immigrazione clandestina.
Nel terzo e nel quarto comma, la Repubblica italiana garantisce a tutti i cittadini stranieri, ai quali siano stati negati i
diritti e le libertà democratiche nei loro paesi, di poter esercitare tali diritti nel territorio dello stato italiano, grazie al
diritto di asilo.
L'ultimo comma prevede che nel nostro paese non sia ammessa l'estradizione dello straniero per reati politici. Lo Stato
italiano rifiuta l'estradizione, cioè il rimandare la persona al paese d’origine, di un cittadino straniero che sia ricercato
per reati politici commessi in opposizione a regimi antidemocratici, nei quali vengono attuate politiche
persecutorie nei confronti dei diritti umani.

\nt{L'italia respinge il nazionalismo, L'Italia è uno stato nazionale, non nazionalista riconosce e difende la propria identità rispetto agli altri stati, ma adotta, atteggiamenti aperti, di collaborazione e di integrazione. La sovranità dello stato Italiano è una condizione che comporta la collaborazione con gli altri stati e anche il riconoscimento dell'autorità istituzionale e sovranazzionale.}


\subsection{Art 11 - Il ripudio della guerra}
\textit{“L'Italia ripudia la guerra come strumento di offesa alla libertà degli altri popoli e come mezzo di risoluzione delle
controversie internazionali; consente, in condizioni di parità con gli altri Stati, alle limitazioni di sovranità necessarie
ad un ordinamento che assicuri la pace e la giustizia fra le Nazioni; promuove e favorisce le organizzazioni
internazionali rivolte a tale scopo.”}

L’articolo 11 oppone un netto rifiuto alla guerra e soprattutto apre un nuovo capitolo, quello di una cultura di
pace da diffondere e costruire anche attraverso organismi internazionali.
Proprio nel periodo in cui è stato elaborato il testo costituzionale, l’Organizzazione delle nazioni unite (Onu) stava
muovendo i suoi primi passi: l’essenza dell’articolo 11 riflette proprio la speranza che l’Italia potesse essere inclusa tra i
paesi “amanti della pace”, così come previsto dallo statuto dell’Onu, e potesse essere ammessa a far parte
dell’Organizzazione stessa, obiettivo che è stato raggiunto nel 1955.
L’Italia, pertanto, in base all’articolo 11 e come membro dell’Onu, non ricorrerà alle armi per risolvere
eventuali contrasti con gli altri stati e non invaderà mai il territorio altrui, violando la libertà di altri popoli.
Quindi l’unico tipo di guerra ammesso è la legittima difesa per respingere un attacco armato che minacci
l’esistenza e l’indipendenza dell’Italia.

\subsection{Art 12 - La bandiera italiana}

\textit{“La bandiera della Repubblica è il tricolore italiano: verde, bianco e rosso, a tre bande verticali di eguali dimensioni.”}

Il tricolore italiano venne decretato il 7 gennaio del 1797 a Reggio Emilia come bandiera della Repubblica Cispadana.
Durante il regno d’Italia fu aggiunto al tricolore lo stemma sabaudo poi eliminato con l’istituzione della repubblica.
Viene ripresa perché era la bandiera del regno d’Italia e venne ammainata definitivamente il 25 aprile 1945, con lo
scioglimento dal giuramento per militari e civili, quale ultimo atto del governo di Benito Mussolini. Il tricolore
rappresenta l’Italia in qualsiasi manifestazione e anche all’estero.

\end{document}
